% Shared preamble: fonts, colours, code listings, tables.
\usepackage[T1]{fontenc}
\usepackage[utf8]{inputenc}
\usepackage{lmodern}
\usepackage{microtype}
\usepackage{amsmath}
\usepackage{graphicx}
\usepackage{booktabs}
\usepackage{tabularx}
\usepackage{longtable}
\usepackage{enumitem}
\usepackage{float}
\usepackage[table]{xcolor}
\usepackage{listings}
\usepackage{caption}
\usepackage{titlesec}
\usepackage{fancyhdr}
\usepackage{xurl}
\usepackage[hidelinks]{hyperref}

\graphicspath{{../figures/}{../../participant-kit/figures/}{../../data/}}

% Colours lifted from participant-kit/plotstyle.py, so the document and the
% charts in it are the same palette.
\definecolor{ink}{HTML}{0B0B0B}
\definecolor{inksoft}{HTML}{52514E}
\definecolor{inkmuted}{HTML}{8A8983}
\definecolor{accent}{HTML}{2A78D6}
\definecolor{warn}{HTML}{EC835A}
\definecolor{good}{HTML}{0CA30C}
\definecolor{rule}{HTML}{E6E5E1}
\definecolor{surface}{HTML}{FCFCFB}

\hypersetup{colorlinks=true, linkcolor=accent, urlcolor=accent}

\captionsetup{font=small, labelfont={bf,color=ink},
              textfont={color=inksoft}, skip=6pt}
\captionsetup[table]{skip=4pt}

\titleformat{\section}{\Large\bfseries\color{ink}}{\thesection}{0.7em}{}
\titleformat{\subsection}{\large\bfseries\color{ink}}{\thesubsection}{0.7em}{}
\titleformat{\subsubsection}{\normalsize\bfseries\color{inksoft}}%
            {\thesubsubsection}{0.7em}{}
\titlespacing*{\section}{0pt}{2.2ex plus 1ex minus .2ex}{1.2ex plus .2ex}

\pagestyle{fancy}
\fancyhf{}
\fancyhead[L]{\small\color{inkmuted}TYTFS 2024 as a PyPSA network}
\fancyfoot[C]{\small\color{inkmuted}\thepage}
\renewcommand{\headrulewidth}{0.4pt}
\renewcommand{\footrulewidth}{0pt}

\lstdefinestyle{code}{
  basicstyle=\ttfamily\scriptsize\color{ink},
  keywordstyle=\color{accent}\bfseries,
  commentstyle=\color{inkmuted}\itshape,
  stringstyle=\color{good},
  numberstyle=\ttfamily\tiny\color{inkmuted},
  backgroundcolor=\color{surface},
  frame=leftline, framerule=1.6pt, rulecolor=\color{rule},
  framexleftmargin=6pt, xleftmargin=10pt,
  breaklines=true, breakatwhitespace=true,
  showstringspaces=false, tabsize=4,
  aboveskip=1.1em, belowskip=1.1em,
  columns=fullflexible, keepspaces=true,
}
\lstset{style=code, language=Python}
\lstdefinelanguage{shell}{morekeywords={python,pip,pytest,git},
  sensitive=true, morecomment=[l]{\#}}

% \code{...} for an inline identifier, \file{...} for a path.
\newcommand{\code}[1]{\texttt{\small #1}}
\newcommand{\file}[1]{\texttt{\small\color{inksoft} #1}}
% \where{file}{symbol} - the standard "this lives here" pointer.
% A long identifier list has no break points once it is in \texttt, so the
% box is set ragged-right with plenty of stretch rather than overfull.
\newcommand{\where}[2]{\par\smallskip\noindent%
  \begingroup\small\color{inksoft}\raggedright\emergencystretch=3em%
  \textbf{Code:}~\texttt{#1}~$\rightarrow$~\texttt{#2}\par\endgroup\smallskip}

\setlist[itemize]{leftmargin=1.3em, itemsep=2pt, topsep=4pt}
\setlist[enumerate]{leftmargin=1.6em, itemsep=2pt, topsep=4pt}
\renewcommand{\arraystretch}{1.15}
